\documentclass[../../../include/open-logic-section]{subfiles}

\begin{document}

\olfileid{sth}{cardinals}{hp}
\olsection{ضميمه: د هيوم اصل}

په \olref[cp]{sec} کښې مو د کانتور اصل بيان کړ۔ هغه دا و:
\begin{align*}
	\card{A} = \card{B} & \text{ هغه وخت او يوازې هغه وخت چې } A \approx B.
\intertext{دا له هغه څه سره ډېر ورته دے چې اوس ورته
\emph{د هيوم اصل} ويل کېږي؛ هغه وايي:}
	\fregenum{x} {F(x)} = \fregenum{x}{G(x)} & \text{ هغه وخت او يوازې هغه وخت چې } F \sim G
\end{align*}
دلته «$F \sim G$» په دې معنا لنډه نښه ده چې د $F$ګانو شمېر
عين د $G$ګانو له شمېر سره برابر دے؛ يعنې $F$ګانې له $G$ګانو
سره په يوه دوه اړخيزه يو پر يو سمون کښې اېښودل کېدے شي، يعنې:
\begin{align*}
	\exists R(&\forall v\forall y(Rvy \lif (Fv \land Gy)) \land {}\\
		&\forall v(Fv \lif \lexists![y][Rvy]) \land {}\\
		&\forall y(Gy \lif \lexists![v][Rvy]))
\end{align*}
خو د هيوم اصل او د کانتور اصل تر منځ د ډول توپير شته۔ د کانتور
د اصل په بيان کښې «$A$» او «$B$» د لومړۍ درجې ترمونه دي چې د
\emph{سټونو} استازيتوب کوي۔ د هيوم د اصل په بيان کښې «$F$»،
«$G$» او «$R$» د لومړۍ درجې ترمونه \emph{نه} دي؛ بلکې د
\emph{پريديکات په ځاے} کښې دي۔ (ښايي د \emph{خاصيتونو}
استازيتوب کوي۔) نو د هيوم اصل په نثري ډول داسې بيانولے شو: د
$F$ګانو شمېر د $G$ګانو له شمېر سره هغه وخت او يوازې هغه وخت
برابر دے چې $F$ګانې له~$G$ګانو سره دوه اړخيزه يو پر يو سمون
ولري۔ دې ته \emph{د هيوم اصل} ويل کېږي، ځکه چې هيوم يو وخت
داسې ليکلي وو:
\begin{quote}
  کله چې دوې شمېرې داسې سره يو ځاے شي چې د يوې هر واحد ته
  د بلې يو واحد ځواب وي، نو موږ يې برابرې بولو۔
  \citep[Pt.III Bk.1 \S1]{Hume1740}
\end{quote}
او \citet[\S63]{Frege1884} د هيوم له دې عبارت څخه نقل وکړ،
بيا يې (په اصل کښې) هغه څه تشريح کړل چې موږ د هيوم اصل نومولے
دے؛ له دې سره د هيوم اصل په اوسني رياضيکي منطق کښې مشهور شو۔

د هيوم د اصل او د فرېګه د پنځم بنسټيز قانون تر منځ جوړښتي ورته
والي ته پام وکړئ۔ موږ هغه په \olref[story][blv]{sec} کښې داسې
بيان کړے و:
\[
	\fregeext{x}{F(x)} = \fregeext{x}{G(x)} \text{هغه وخت او يوازې هغه وخت چې } \lforall[x][(F(x) \liff G(x))].
\]
دلته يو څه تشريح او پرتله ګټوره ده۔

د هيوم د اصل يا د پنځم بنسټيز قانون په شان اصل په دوو ډولونو
اخيستلے شو: \emph{پريديکاتيف} يا \emph{امپريديکاتيف} (وګورئ
\olref[story][predicative]{sec})۔ د پنځم بنسټيز قانون په
امپريديکاتيف لوست کښې، د هر~$F$ دپاره څيز
$\fregeext{x}{F(x)}$ د کمي کولو په هماغه دامن کښې راځي چې د
پنځم بنسټيز قانون د بيان دپاره مو کارولې وه۔ په ورته ډول، د
هيوم د اصل په امپريديکاتيف لوست کښې، د هر~$F$ دپاره څيز
$\fregenum{x}{F(x)}$ د کمي کولو په هماغه دامن کښې راځي چې د
هيوم د اصل د بيان دپاره مو کارولې وه۔ برعکس، په
\emph{پريديکاتيف} لوست کښې به څيزونه
$\fregeext{x}{F(x)}$ او~$\fregenum{x}{F(x)}$ د کوم
\emph{بل} دامن موجودات وي۔

که پنځم بنسټيز قانون امپريديکاتيف ولولو، د ساده ادراک له لارې
تناقض ته رسېږي (د تفصيل دپاره \olref[story][blv]{sec} وګورئ)۔
د ساده ادراک په څېر، د هغه پريديکاتيف لوست يې له تناقضه
ساتلے شي۔ خو ښايي بيا هغه ټول کارونه ونه کړي چې موږ ترې غوښتل۔

خو د هيوم اصل \emph{کېدے شي} په امپريديکاتيف ډول بې له تناقضه
ولوستل شي۔ او په دې لوست کښې ډېر ځواکمن دے۔

د روښانولو دپاره پريديکات «$x \neq x$» وګورئ؛ څرګنده ده چې هېڅ
څيز يې نه پوره کوي۔ د هيوم اصل اوس يو څيز $\# x( x\neq x)$
راکوي۔ دا ښايي د~$0$ عدد وګڼو۔ اوس د \emph{امپريديکاتيف} لوست
له مخې---خو \emph{يوازې} له همدې لوست سره---دا څيز $0$ زموږ
د کمي کولو په اصلي دامن کښې شامل دے۔ نو د هيوم اصل پر پريديکات
«$x = 0$» په معنا لرونکي ډول کارولے شو او څيز
$\#x (x = 0)$ ترلاسه کولے شو۔ دا ښايي د~$1$ عدد وګڼو۔ سربېره
پر دې، د هيوم اصل وايي چې $0 \neq 1$، ځکه د هغو څيزونو چې له
ځان سره برابر نه دي او هغو چې له $0$ سره برابر دي تر منځ دوه
اړخيز يو پر يو سمون نشته (له لومړيو هېڅ نشته، خو له دويمو يو
شته)۔ اوس بيا په امپريديکاتيف لوست کښې $1$ زموږ د کمي کولو په
اصلي دامن کښې شامل دے۔ نو د هيوم اصل پر پريديکات
«$(x = 0 \lor x = 1)$» کارولے شو او څيز
$\#x(x = 0 \lor x = 1)$ ترلاسه کولے شو۔ دا ښايي د~$2$ عدد
وګڼو، او ښودلے شو چې $0\neq 2$ او $1 \neq 2$، او همداسې نور۔

په لنډه، د هيوم اصل په امپريديکاتيف لوست کښې د \emph{نامتناهي
ډېرو څيزونو} شتون راکوي۔ دې خبرې \emph{نويو فرېګه‌پالو
منطق‌پوهانو} ته هڅونه ورکړې چې د هيوم اصل د حساب بنسټ ونيسي۔

خو فرېګه \emph{پخپله} د هيوم اصل د حساب بنسټ نۀ ګرځاوه۔ د دې
پر ځاے يې د هيوم اصل له يوه څرګند تعريفه ثابت کړ:
$\fregenum{x}{F(x)}$ د مفهوم $F \sim \Phi$ د غځونې په توګه
تعريفېږي۔ په اوسنۍ ژبه کښې ښايي دا په
$\fregenum{x}{F(x)} = \Setabs{G}{F \sim G}$ بيان کړو؛ خو دا مو
بيا د ساده ادراک ستونزو ته راګرځوي۔

\end{document}
